
Three uncertainty quantification (UQ) signals — token-level entropy, semantic entropy, and SelfCheckGPT-BERTScore — have been proposed for zero-resource hallucination detection in large language models, but have been validated in separate settings on distinct benchmarks, making cross-method comparisons difficult. This paper presents a controlled comparison of all three methods applied to the same model (LLaMA-2-7B-chat) on a stratified 2,000-example subset of HaluEval-QA under matched inference conditions. All three methods fail to discriminate hallucinated from factual responses: AUROC values are 0.5000 (semantic entropy), 0.4829 (token entropy), and 0.3562 (SelfCheckGPT-BERTScore). The failure of semantic entropy is traceable to a specific mechanism: the NLI clustering step (microsoft/deberta-large-mnli) assigns all five stochastic responses to distinct semantic clusters in 72.8\% of examples, yielding an aggregation rate of 0.272 (95\% CI [0.253, 0.292]) and collapsing semantic entropy to a constant signal. SelfCheckGPT-BERTScore achieves below-random discrimination (AUROC = 0.3562, 95\% CI [0.332, 0.380]), consistent with a label polarity inversion hypothesis — that ChatGPT-generated confident hallucinations in HaluEval-QA are produced consistently by LLaMA-2-7B-chat, inverting the consistency–hallucination correlation that SelfCheckGPT depends on. This hypothesis is not verified in the present work. The cross-model comparison (Mistral-7B-Instruct) was not executed. The findings quantify domain boundaries for NLI-based UQ methods on short factual QA responses.

